\section{Comparación de métodos: simplificar el proceso no significa dejar de tener supuestos}
El atractivo de DPO radica en que sustituye el proceso de dos etapas de RLHF por una cadena de optimización más corta, pero esto no significa que carezca de premisas. Es válido porque aceptamos una relación específica entre la política óptima y el reward, y porque estamos dispuestos a usar la reference policy y la restricción de KL para mantener estable el entrenamiento.

\begin{knowledgebox}{Tres puntos que hay que conservar al entender DPO}
\begin{itemize}
  \item Lo que elimina es el RM explícito, no el modelado de preferencias en sí
  \item La reference policy sigue desempeñando el papel de estabilizador en el objetivo
  \item La calidad de los datos de preferencias sigue determinando de manera directa el límite superior del entrenamiento
\end{itemize}
\end{knowledgebox}

\subsection{Resumen de la sección}
La «elegancia» de DPO proviene de la sencillez de su derivación, pero en la práctica sigue siendo necesario atender la calidad de los datos, la elección de la reference y el ruido en las preferencias; no debe entenderse como un sustituto sin costo.

\newpage
